% Generated draft: requires mathematical and source review.
\documentclass[11pt]{article}
\usepackage[margin=25mm]{geometry}
\usepackage{fontspec}
\setmainfont{Latin Modern Roman}
\usepackage{amsmath,amssymb,mathtools,unicode-math}
\setmathfont{Latin Modern Math}
\usepackage{enumitem,needspace,xurl,hyperref}
\setlength{\parindent}{0pt}
\setlength{\parskip}{5pt}
\setlength{\emergencystretch}{3em}
\newcommand{\N}{\mathbb{N}}
\newcommand{\Z}{\mathbb{Z}}
\newcommand{\Q}{\mathbb{Q}}
\newcommand{\R}{\mathbb{R}}
\newcommand{\C}{\mathbb{C}}
\newcommand{\F}{\mathbb{F}}
\newcommand{\A}{\mathbb{A}}
\newcommand{\PP}{\mathbb{P}}
\newcommand{\E}{\mathbb{E}}
\newcommand{\eps}{\varepsilon}
\newcommand{\dd}{\,\mathrm{d}}
\begin{document}
\section*{UnsolvedMath ID: 700776}
\textbf{Source problem code:} catalogue entry
\subsection{Definitions and mathematical statement}
For a group \(G\), define its derived series by \(G^{(0)}=G\) and \(G^{(i+1)}=[G^{(i)},G^{(i)}]\), where \([H,H]\) is the subgroup generated by commutators of elements of \(H\). The group \(G\) is \emph{metabelian} if \(G^{(2)}=\{1\}\), and \emph{soluble} if \(G^{(d)}=\{1\}\) for some integer \(d\geq 0\). It is \emph{finitely generated} if some finite subset generates it.

Under the usual interpretation of \emph{finite rank} as finite Pr\"ufer rank, \(G\) has finite rank if there exists an integer \(r\geq 0\) such that every finitely generated subgroup of \(G\) can be generated by at most \(r\) elements. The supplied record does not explicitly define rank, so this interpretation requires confirmation.

For a class \(\mathcal C\) of groups with a specified effective input representation, its \emph{isomorphism problem} asks whether there exists an algorithm which, for every pair of valid descriptions representing groups \(G,H\in\mathcal C\), halts and determines whether \(G\cong H\). Here \(G\cong H\) means that there exists a bijection \(\varphi:G\to H\) satisfying \(\varphi(xy)=\varphi(x)\varphi(y)\) for all \(x,y\in G\). No behavior on descriptions outside the promised class is required by this formulation.

\textbf{Problem (V. N. Remeslennikov, 1990; Kourovka Notebook 11.84).} Is the isomorphism problem decidable separately for each of the following classes?
\begin{enumerate}
\item[(a)] All finitely generated metabelian groups.
\item[(b)] All finitely generated soluble groups of finite rank.
\end{enumerate}
The intended algorithm in each part must work uniformly over the stated class. No common bound on the number of generators, derived length, or finite rank is supplied.

\textbf{Unresolved input qualification.} The supplied statement does not specify how the groups are given effectively, or whether bounds or other auxiliary data accompany their descriptions. Finite generation alone does not specify an algorithmic input format. Consequently, the questions above preserve the source's intended decision problems but require an input convention before they become fully specified algorithmic statements; in particular, no restriction to groups having finite ordinary presentations has been imposed.
\subsection{Short English statement}
Remeslennikov asks whether an algorithm can decide when two finitely generated metabelian groups are isomorphic, and separately whether this is possible for finitely generated soluble groups of finite rank. The source leaves the effective descriptions of the groups unspecified, and the meaning of finite rank requires confirmation. Its supplied status is \emph{open}; this draft does not independently verify that status.
\subsection{Sources}
\begin{itemize}
\item Local UnsolvedMath registry, supplied complete source record: permanent ID \(700776\), exact source code \texttt{catalogue entry}, title \emph{Kourovka Notebook Problem 11.84}. The supplied record labels the problem \emph{open}.
\item \emph{The Kourovka Notebook (Unsolved Problems in Group Theory)}, original problem 11.84, 11th issue (1990); proposed by V. N. Remeslennikov. The supplied background also identifies the 21st issue edition (2026), arXiv identifier \texttt{1401.0300v48}. The supplied collection description names E. I. Khukhro and V. D. Mazurov as editors.
\item Registry provenance: upstream revision \texttt{d371ee917ac9647fbb06cd6fc529deec6f3dda82}; record created and updated at \texttt{2026-10-16T14:00:00Z}. These are supplied registry timestamps, not independent literature-verification dates. No separate literature-check date or literature-source links are supplied. The cited Notebook edition and the registry's status assessment have not been independently checked for this draft.
\end{itemize}
\end{document}
